\documentclass[12pt]{article}
\usepackage{amsmath}
\usepackage{amssymb}
\usepackage{geometry}
\geometry{a4paper, margin=1in}
\usepackage{graphicx}

\title{Traslaciones y Representación de Objetos mediante Matrices en Juegos 2D}
\author{Curso de Desarrollo de Videojuegos 2D}
\date{}

\begin{document}

\maketitle

\section*{Introducción}
En juegos 2D, los objetos se representan mediante vectores de posición y se pueden manipular mediante operaciones básicas de álgebra lineal. En esta sección, describimos cómo se modelan los objetos como matrices de vectores y cómo se aplican traslaciones simples para el movimiento dentro del juego.

\section*{Representación de un objeto}
Un objeto 2D puede representarse mediante un vector de posición:
$$
[
\mathbf{v} = \begin{bmatrix} x \ y \end{bmatrix},]
$$
donde $x$ e $y$ representan la coordenada del objeto en la pantalla.

Si el objeto tiene múltiples puntos de interés (como vértices de un rectángulo), se puede representar como una matriz cuyos columnas son los vectores de cada punto:
$$
[
\mathbf{O} = \begin{bmatrix} x_1 & x_2 & x_3 & x_4 \\ y_1 & y_2 & y_3 & y_4 \end{bmatrix}]
$$
para un objeto rectangular de 4 vértices.

\section*{Traslaciones}
La traslación es una operación que desplaza un objeto por un vector $(dx, dy)$. Para un vector de posición:
$$
[
\mathbf{v}' = \mathbf{v} + \mathbf{t} = \begin{bmatrix} x \ y \end{bmatrix} + \begin{bmatrix} dx \ dy \end{bmatrix} = \begin{bmatrix} x + dx \ y + dy \end{bmatrix}.]
$$

Para un objeto con varios puntos representado por la matriz $\mathbf{O}$:
$$
[
\mathbf{O}' = \mathbf{O} + \mathbf{T},]
$$
donde $\mathbf{T}$ es una matriz con el vector de traslación repetido en cada columna:
$$
[
\mathbf{T} = \begin{bmatrix} dx & dx & dx & dx \\ dy & dy & dy & dy \end{bmatrix}.]
$$

\section*{Aplicaciones en el juego}
En nuestro juego 2D de plataformas, las operaciones básicas aplicadas son:
\begin{itemize}
\item \textbf{Movimiento horizontal del jugador:} Desplazamiento en $x$ mediante $dx = \pm v$ según la dirección.
\item \textbf{Salto y caída:} Desplazamiento en $y$ con $dy$ determinado por la velocidad vertical y gravedad.
\item \textbf{Movimiento de enemigos o cámara:} Traslación de todos los objetos visibles en sentido contrario al movimiento del jugador para simular seguimiento de cámara.
\end{itemize}

\section*{Resumen}
Cada objeto del juego puede conceptualizarse como un conjunto de vectores. Las traslaciones se implementan mediante suma de vectores o matrices de vectores, y constituyen la operación básica para mover personajes, enemigos, objetos y la cámara dentro del juego 2D.

\end{document}
